\documentclass[11pt]{article}
\usepackage{amsmath,amssymb,amsthm}
\usepackage[margin=1in]{geometry}
\usepackage{hyperref}
\newtheorem{lemma}{Lemma}
\newtheorem{theorem}{Theorem}
\newtheorem{corollary}{Corollary}
\newtheorem{proposition}{Proposition}
\newcommand{\ii}{\mathrm i}
\title{A two-sided certificate for continuous-domain scalar diffraction residuals}
\author{}\date{}
\begin{document}\maketitle
\paragraph{Definitions and notation.}
Let $A,\Omega\subset\mathbb R^2$ be bounded measurable sets of positive
area, with probability measure $d\nu=d\mathbf r/|\Omega|$ on $\Omega$.
Fix real $\lambda,z>0$, $K=2\pi z/\lambda$ and
$q=|\mathbf r-\mathbf s|^2/z^2\in[0,Q]$, $0<Q<\infty$. Let $t=\sqrt{1+q}$
be the principal branch. All three fields use unit incident amplitude,
the common multiplier $e^{\ii K}/(\ii\lambda z)$ and kernels
\[
 H_R=a e^{\ii K\phi},\quad a=(1+q)^{-1}+\ii/(Kt^3),\quad
 \phi=q/(1+t),\qquad H_F=e^{\ii Kq/2},
\]
\[
 H_P=b e^{\ii K\psi},\quad
 b=(1+q)^{-1}+\frac{\ii(1+q/4-q^2/48)}{K(1+7q/4+35q^2/48)},\quad
 \psi=\frac{q/2+q^2/4}{1+3q/4+q^2/16}.
\]
Choose $D=H_R-H_P$ or $H_R-H_F$. Removing the common unit-modulus
carrier leaves the absolute norm invariant. No fit of phase or amplitude is made.
Write $Ef=(\ii\lambda z)^{-1}\int_A f(q(\mathbf r,\mathbf s))\,d\mathbf s$,
$N=\|ED\|_{L^2(\nu)}$.

\begin{lemma}[Uniform spectral interpolation error]
Suppose $f(\xi)=D(Q(1+\xi)/2)$ is holomorphic in a neighborhood of
the filled closed Bernstein ellipse of parameter $\rho>1$ and $|f|\le M$
there. For integer $n\ge1$, let $P_n$ interpolate $D$ at
$Q(1+\cos(j\pi/n))/2$, $0\le j\le n$. Then
$\|D-P_n\|_{\infty,[0,Q]}\le\epsilon_n=4M\rho^{-n}/(\rho-1)$.
\end{lemma}
\begin{proof}
The Laurent expansion after $\xi=(w+w^{-1})/2$ gives
$|a_j|\le2M\rho^{-j}$ for positive Chebyshev indices by Cauchy's formula.
On the Lobatto nodes $T_j$ aliases to $T_{\ell}$ where $\ell$ is obtained
by reducing modulo $2n$ and reflecting into $[0,n]$. Uniform absolute
convergence permits termwise interpolation. Each tail term contributes at
most $2|a_j|$ since $|T_j-T_\ell|\le2$ on $[-1,1]$. The geometric
sum gives the claim. This argument applies to complex coefficients
directly. The classical truncation estimate has a different constant
\cite[Thm.~8.2]{tref}.
\end{proof}

\begin{lemma}[Exact moment representation]
For a polynomial $P_n$, $EP_n$ is a polynomial in the observer coordinates.
Given its coefficients $C_\alpha$ in a real monomial basis $b_\alpha$,
the squared norm is $\sum_{\alpha,\beta} C_\alpha\overline{C_\beta}
\int_\Omega b_\alpha b_\beta\,d\nu$.
\end{lemma}
\begin{proof}
Expand $((x-s_x)^2+(y-s_y)^2)^d$ by the binomial theorem. Each
source monomial is integrable by boundedness of $A$; the sums are
finite. The Hermitian identity follows by pointwise multiplication and
integration. Computing $(\Re EP_n)^2+(\Im EP_n)^2$ gives the same
signed cross terms. A global unit phase drops out of this expression.
\end{proof}

\begin{theorem}[Certified continuous-domain $L^2$ error bounds]
Under the assumptions of Lemma~1 set $G_n=\|EP_n\|_{L^2(\nu)}^2$
and $\delta_n=|A|\epsilon_n/(\lambda z)$. Then
\[
 \max(0,\sqrt{G_n}-\delta_n)\ \le N\le\sqrt{G_n}+\delta_n .
\]
\end{theorem}
\begin{proof}
For every observer, the magnitude of $E(D-P_n)$ is bounded by
$\int_A\epsilon_n/(\lambda z)=\delta_n$. Because $\nu(\Omega)=1$,
the same constant bounds its $L^2(\nu)$ norm. Apply the reverse triangle
inequality to $ED,EP_n$, and use $N\ge0$. This estimate leaves all
phase cancellation inside $EP_n$ intact. Lebesgue area in place of
$\nu$ would produce an additional factor $\sqrt{|\Omega|}$ in the slack.
For prescribed nonuniform illumination $w$, replace $|A|$ by
$\int_A |w|$ and use source moments weighted by $w$.
\end{proof}

\begin{corollary}[Superelliptical apertures]
Let $A$ be a centered axis-aligned superellipse with positive semiaxes
$a_x,a_y$ and $p>0$. Its even moments are
\[
 \mu_{ij}=\frac{4a_x^{2i+1}a_y^{2j+1}
 \Gamma((2i+1)/p)\Gamma((2j+1)/p)}
 {p^2\Gamma(1+(2i+2j+2)/p)}.
\]
For a concentric disk of radius $R$, take $Q\ge (R+r_A)^2/z^2$ with
\[
 r_A=\begin{cases}\max(a_x,a_y),&0<p\le2,\\
 (a_x^{2p/(p-2)}+a_y^{2p/(p-2)})^{(p-2)/(2p)},&p>2.
 \end{cases}
\]
With $X=(x/R)^2$, $Y=(y/R)^2$, the disk average of $X^iY^j$ is
$(1/2)_i(1/2)_j/(i+j+1)!$.
\end{corollary}
\begin{proof}
In each quadrant substitute $u=(|s_x|/a_x)^p$,
$v=(|s_y|/a_y)^p$. Integrating $u^{(2i+1)/p-1}v^{(2j+1)/p-1}$
over $u+v\le1$ gives the Dirichlet integral above. Reflection removes
odd moments. The area is $\mu_{00}$ and $AR=a_x/a_y$ fixes the axis ratio.
For $p\le2$, $u^{2/p}+v^{2/p}\le u+v\le1$; the longer axis attains
the bound. For $p>2$, apply H\"older with conjugate exponents
$p/(p-2),p/2$ to $(a_x^2,a_y^2)$ and $(u^{2/p},v^{2/p})$.
The equality case is attained on the source boundary. An opposite
point on the observer rim realizes the maximum separation.
For the disk moment, polar integration gives a radial factor
$1/(2i+2j+2)$ and an angular integral
$2B(i+1/2,j+1/2)$, divided by $\pi$ after rescaling to unit radius.
\end{proof}

\begin{proposition}[Validated interval-arithmetic certificate]
Assume verified finite enclosures $G_n\in[L,H]$, $0\le\delta_n\le d$,
$L\le H$, $H\ge0$. Then valid endpoints are the outward roundings of
$\max(0,\sqrt{\max(0,L)}-d)$ and $\sqrt H+d$.
\end{proposition}
\begin{proof}
Monotonicity of the square root on $[0,\infty)$ and of subtraction
in its first argument gives the result. $H<0$ is incompatible with
the squared norm and must abort certification. A negative $L$ is
consistent with an enclosure of a nonnegative exact quantity.
An implementation must also establish validity of its physical domain,
model identity, contour covering and moment enclosures before using
the proposition.
\end{proof}

\begin{corollary}[Geometric convergence]
For fixed admissible $\rho,M,A,\lambda,z$ the exact-arithmetic
width is at most $8|A|M\rho^{-n}/((\rho-1)\lambda z)$.
\end{corollary}
\begin{proof}
The width is at most $2\delta_n$. A rounded squared-norm enclosure
can add $\sqrt{H}-\sqrt{\max(0,L)}$ plus endpoint rounding. Fixed
working precision alone does not imply geometric convergence.
\end{proof}
\paragraph{Assumptions and limitations.}
The general theorem does not require reflection symmetry or a disk.
Those assumptions enter the particular moment evaluator. For the two
specified scalar kernels the holomorphic neighborhood must be checked
for each geometry and parameter choice. The construction combines
classical interpolation, moment identities and norm inequalities.
No assertion of priority for these results, a Maxwell boundary model,
or universal certification by an unguarded implementation follows.
\begin{thebibliography}{1}
\bibitem{tref} L.\ N.\ Trefethen, \emph{Approximation Theory and
Approximation Practice}, SIAM, 2013, Theorems 8.1--8.2.
\url{https://people.maths.ox.ac.uk/trefethen/trefethen_sample.pdf}.
\end{thebibliography}
\end{document}
